% Compile from repository root with:
% tectonic professor-report-final.tex
\documentclass[11pt,a4paper]{article}

\usepackage[utf8]{inputenc}
\usepackage[T1]{fontenc}
\usepackage[italian]{babel}
\usepackage{geometry}
\usepackage{graphicx}
\usepackage{booktabs}
\usepackage{array}
\usepackage{longtable}
\usepackage{float}
\usepackage{xcolor}
\usepackage{amsmath}
\usepackage{xurl}
\usepackage{hyperref}

\geometry{margin=2.1cm}
\hypersetup{
  colorlinks=true,
  linkcolor=blue,
  urlcolor=blue,
  citecolor=blue
}

\newcommand{\normal}{\textit{Normal}}
\newcommand{\pneumonia}{\textit{Pneumonia}}
\newcommand{\balacc}{balanced accuracy}
\newcommand{\auc}{ROC-AUC}
\newcommand{\githuburl}{\url{https://github.com/yahiag04/pneumonia-xray-classifier}}

\title{\textbf{Classificazione binaria di polmonite da radiografie toraciche -- Inquadramento per tesi}\\
\large Confronto tra reti CNN, generalizzazione cross-dataset e analisi statistica}
\author{Yahia Ghallale\\
\small Repository GitHub: \githuburl}
\date{Università degli Studi di Brescia \\ 1 luglio 2026}

\begin{document}
\maketitle

\begin{abstract}
Questo documento riassume lo stato finale degli esperimenti svolti finora per
la tesi sulla classificazione binaria \normal/\pneumonia{} da radiografie
toraciche. Il lavoro e' stato riorientato verso un setting adulto: il training
principale e' stato effettuato su un sottoinsieme binario e bilanciato del
dataset RSNA Pneumonia Detection Challenge, mentre le prestazioni sono state
misurate su test interno RSNA e su due dataset esterni, Kermany e Chittagong.
Sono state confrontate cinque reti: una CNN custom, denominata PneumoniaNet, e
quattro architetture note, ResNet18, MobileNetV3-Large, EfficientNet-B0 e
DenseNet121. In aggiunta sono stati eseguiti test zero-shot con modelli
TorchXRayVision pre-addestrati su dataset radiografici. L'analisi finale mostra
che non esiste una rete migliore in modo uniforme su tutti i dataset:
EfficientNet-B0 e' la migliore sul test interno RSNA, DenseNet121 e' la piu'
forte su Kermany, mentre su Chittagong PneumoniaNet rientra nel gruppo di testa
ma e' statisticamente indistinguibile da ResNet18, MobileNetV3-Large ed
EfficientNet-B0. L'analisi delle soglie mostra inoltre che il ranking su
Chittagong e' molto sensibile al punto operativo, mentre il gap
Kermany--Chittagong dei singoli modelli segue da vicino la differenza di AUC.
L'analisi di calibrazione mostra che Chittagong ha errore di calibrazione piu'
alto di RSNA per tutte le architetture, ma una calibrazione post-hoc tramite
temperature scaling fittata su validation RSNA non basta a stabilizzare il
ranking operativo su Chittagong. Nel complesso, la selezione del modello in
trasferimento cross-domain richiede dati target-domain, o almeno una validation
rappresentativa del dominio di deploy. E' stata inoltre aggiunta una prima
estensione esplorativa multi-classe RSNA, in cui la classe esclusa
\texttt{No Lung Opacity / Not Normal} viene trattata come terza categoria
distinta.
\end{abstract}

\tableofcontents
\newpage

\section{Obiettivo}

L'obiettivo della tesi e' valutare modelli di deep learning per la
classificazione binaria di radiografie toraciche in due classi:
\normal{} e \pneumonia. Il punto centrale non e' soltanto ottenere una buona
accuracy su un dataset, ma capire quanto le reti generalizzino quando vengono
testate su distribuzioni diverse.

La domanda sperimentale finale e':

\begin{quote}
\textit{Una rete addestrata in modo controllato su un dataset adulto e
pneumonia-specifico mantiene prestazioni robuste su dataset esterni? E una
singola architettura risulta migliore in modo stabile su tutti i domini?}
\end{quote}

Per rispondere sono stati separati due tipi di esperimento:

\begin{itemize}
  \item \textbf{zero-shot}: modelli gia' pre-addestrati su radiografie
  toraciche vengono valutati senza ulteriore training sui dataset esterni;
  \item \textbf{post-training}: le reti del progetto vengono addestrate sul
  training set RSNA adult-oriented e poi valutate su RSNA, Kermany e Chittagong.
\end{itemize}

\section{Dataset}

\subsection{RSNA adult-oriented}

Il dataset principale e' RSNA Pneumonia Detection Challenge
\cite{rsna_kaggle}. La scelta e' motivata dal fatto che il task e'
pneumonia-specifico e che la classe positiva deriva da annotazioni radiologiche
di opacita' polmonare, non da una semplice binarizzazione di un dataset
multi-label. La versione locale usata negli esperimenti contiene immagini PNG
processate e metadata compatibili con la preparazione binaria.

Per ottenere un task coerente con \normal/\pneumonia{} e' stata usata la
seguente mappatura:

\begin{itemize}
  \item \texttt{Normal} $\rightarrow$ \normal;
  \item \texttt{Lung Opacity} $\rightarrow$ \pneumonia;
  \item \texttt{No Lung Opacity / Not Normal} esclusa dal task binario.
\end{itemize}

E' stato costruito un subset bilanciato da 7000 immagini: 5000 per training,
1000 per validation e 1000 per test. Lo split e' stato fatto a livello di
\texttt{patientId}: prima dello split e' stata mantenuta una sola immagine per
paziente e poi i pazienti sono stati divisi tra training, validation e test.
Nel metadata finale ci sono 7000 righe e 7000 \texttt{patientId} unici; le
intersezioni train/validation, train/test e validation/test sono tutte nulle.
Questo evita leakage interno tra gli split RSNA.

I metadata RSNA includono anche la proiezione radiografica. Nel subset finale
sono presenti 3635 immagini PA e 3365 AP. La distribuzione per split e':
training 2621 PA e 2379 AP, validation 493 PA e 507 AP, test 521 PA e 479 AP.
I dataset esterni non hanno nel manifest locale lo stesso livello di metadata
DICOM; la view position rimane quindi un possibile confondente non completamente
controllabile nella valutazione cross-dataset.

Come estensione esplorativa e' stato poi costruito anche un task RSNA
multi-classe a tre classi, mantenendo la classe precedentemente esclusa:

\begin{itemize}
  \item \texttt{Normal} $\rightarrow$ \texttt{normal};
  \item \texttt{Lung Opacity} $\rightarrow$ \texttt{lung\_opacity};
  \item \texttt{No Lung Opacity / Not Normal} $\rightarrow$
  \texttt{not\_normal\_no\_lung\_opacity}.
\end{itemize}

Lo split multi-classe e' bilanciato a livello di classe: 2500 immagini per
classe in training, 500 per classe in validation e 500 per classe in test. Il
test multi-classe contiene quindi 1500 immagini totali, equamente divise tra le
tre categorie. Questa tabella descrive lo split disponibile; l'esperimento
multi-classe riportato piu' avanti non usa ancora il training set come
fine-tuning end-to-end completo, ma come linear probe sui checkpoint binari.

\subsection{Dataset esterni}

Sono stati usati due dataset esterni:

\begin{itemize}
  \item \textbf{Kermany}: dataset collegato allo studio di Kermany et al.
  \cite{kermany2018}. E' utile come test fuori distribuzione, ma non viene
  trattato come training principale perche' nasce da una popolazione pediatrica
  diversa dal target adult-oriented della tesi.
  \item \textbf{Chittagong}: dataset esterno indipendente con classi
  \normal{} e \pneumonia. E' stato usato solo lo split di test, senza
  fine-tuning, per misurare la generalizzazione esterna.
\end{itemize}

\begin{table}[H]
\centering
\caption{Dataset usati negli esperimenti finali.}
\begin{tabular}{lrrr}
\toprule
Dataset & Normal & Pneumonia & Totale \\
\midrule
RSNA train+val+test & 3500 & 3500 & 7000 \\
RSNA test & 500 & 500 & 1000 \\
Kermany test & 234 & 390 & 624 \\
Chittagong test & 257 & 256 & 513 \\
\bottomrule
\end{tabular}
\end{table}

\begin{table}[H]
\centering
\caption{Split RSNA multi-classe usato nell'estensione esplorativa.}
\begin{tabular}{lrrrr}
\toprule
Split & Normal & Lung opacity & Not normal/no opacity & Totale \\
\midrule
Train & 2500 & 2500 & 2500 & 7500 \\
Validation & 500 & 500 & 500 & 1500 \\
Test & 500 & 500 & 500 & 1500 \\
\bottomrule
\end{tabular}
\end{table}

\section{Modelli confrontati}

\subsection{PneumoniaNet}

PneumoniaNet e' la rete custom del progetto. E' una CNN compatta, pensata come
baseline controllata rispetto ad architetture piu' grandi pre-addestrate su
ImageNet. Lavora su input in scala di grigi e contiene circa 300 mila
parametri, quindi e' molto piu' piccola delle reti di confronto.

\begin{table}[H]
\centering
\caption{Struttura di PneumoniaNet.}
\begin{tabular}{lll}
\toprule
Blocco & Operazioni principali & Canali \\
\midrule
Input & radiografia 224x224 in scala di grigi & 1 \\
Blocco 1 & Conv 3x3 + BatchNorm + ReLU + MaxPool & 16 \\
Blocco 2 & Conv 3x3 + BatchNorm + ReLU, due volte + MaxPool & 32 \\
Blocco 3 & Conv 3x3 + BatchNorm + ReLU, due volte + MaxPool & 64 \\
Blocco 4 & Conv 3x3 + BatchNorm + ReLU, due volte + MaxPool & 128 \\
Pooling & Adaptive average pooling & 128 \\
Classifier & Linear 128$\rightarrow$64 + ReLU + Dropout + Linear 64$\rightarrow$1 & 1 logit \\
\bottomrule
\end{tabular}
\end{table}

\subsection{Reti di confronto}

Le reti di confronto sono state scelte per coprire famiglie CNN note e
ampiamente usate:

\begin{itemize}
  \item \textbf{ResNet18}: architettura residuale, basata su connessioni skip
  che facilitano l'ottimizzazione di reti profonde \cite{he2016resnet};
  \item \textbf{MobileNetV3-Large}: architettura efficiente progettata con
  hardware-aware search e blocchi compatti \cite{howard2019mobilenetv3};
  \item \textbf{EfficientNet-B0}: architettura basata su compound scaling, cioe'
  bilanciamento congiunto di profondita', larghezza e risoluzione
  \cite{tan2019efficientnet};
  \item \textbf{DenseNet121}: architettura con connessioni dense tra layer, che
  favoriscono riuso delle feature e propagazione del gradiente
  \cite{huang2017densenet}.
\end{itemize}

\begin{table}[H]
\centering
\caption{Dimensione, costo computazionale e andamento del training su RSNA. Il GMAC e' calcolato per una inferenza con input 224x224, riportato a quattro decimali, e conta le operazioni MAC di layer convoluzionali e lineari, in linea con analisi pratiche costo/accuratezza di reti neurali profonde \cite{canziani2016analysis}.}
\begin{tabular}{llrrrr}
\toprule
Modello & Input & Parametri & GMAC & Best epoch & Epoche eseguite \\
\midrule
PneumoniaNet & 1x224x224 & 300161 & 0.5275 & 5 & 10 \\
ResNet18 & 3x224x224 & 11177025 & 1.8136 & 10 & 15 \\
MobileNetV3-Large & 3x224x224 & 4203313 & 0.2153 & 5 & 10 \\
EfficientNet-B0 & 3x224x224 & 4008829 & 0.3845 & 4 & 9 \\
DenseNet121 & 3x224x224 & 6954881 & 2.8331 & 6 & 11 \\
\bottomrule
\end{tabular}
\end{table}

La tabella separa due aspetti diversi della leggerezza di un modello.
PneumoniaNet e' nettamente la piu' piccola in termini di parametri, quindi ha
un costo di memoria inferiore rispetto alle architetture pre-addestrate. Non e'
pero' la piu' economica in termini di operazioni per inferenza: MobileNetV3-Large
ed EfficientNet-B0 hanno meno GMAC grazie a blocchi progettati esplicitamente
per efficienza computazionale. Di conseguenza la rete custom va descritta come
compatta in memoria, non come architettura computazionalmente ottimizzata.

\section{Metriche}

Per evitare conclusioni fuorvianti in presenza di dataset sbilanciati, la
metrica principale e' la \balacc:

\[
\text{Balanced Accuracy} =
\frac{\text{Sensitivity} + \text{Specificity}}{2}.
\]

Negli output sperimentali sono state inoltre calcolate accuracy, sensitivity,
specificity, ROC-AUC, PR-AUC, F1-score e confusion matrix; le tabelle finali si
concentrano sulle metriche piu' rilevanti per la discussione, cioe' \balacc{},
sensitivity, specificity e ROC-AUC. Per la parte finale sono stati aggiunti:

\begin{itemize}
  \item \textbf{bootstrap confidence intervals}: intervalli al 95\% tramite
  bootstrap stratificato, ispirati al principio bootstrap di Efron
  \cite{efron1979bootstrap};
  \item \textbf{McNemar test}: test paired sulle predizioni corrette/errate di
  due classificatori sullo stesso dataset \cite{mcnemar1947};
  \item \textbf{DeLong test}: confronto statistico tra AUC correlate, quindi
  calcolate sugli stessi campioni \cite{delong1988};
  \item \textbf{Holm-Bonferroni}: correzione per confronti multipli applicata
  separatamente entro ciascun dataset e tipo di test, cosi' da non interpretare
  come stabile un risultato borderline dovuto al numero di confronti
  \cite{holm1979}.
  \item \textbf{calibrazione}: Expected Calibration Error (ECE) con binning
  equal-width ed equal-mass, Brier score e temperature scaling post-hoc
  \cite{guo2017calibration}. L'ECE misura la differenza media pesata tra
  confidenza predetta e frequenza osservata; il Brier score e' una metrica
  binning-free basata sull'errore quadratico delle probabilita'.
\end{itemize}

\section{Esperimenti zero-shot}

Prima del training adult-oriented e' stato eseguito un esperimento economico:
valutare modelli TorchXRayVision gia' pre-addestrati su radiografie toraciche.
TorchXRayVision fornisce dataset, preprocessing e modelli pre-addestrati per
CXR, ed e' adatto come baseline esterna \cite{torchxrayvision}. Sono stati
usati tre DenseNet121:

\begin{itemize}
  \item \texttt{densenet121-res224-rsna};
  \item \texttt{densenet121-res224-nih};
  \item \texttt{densenet121-res224-all}.
\end{itemize}

\begin{table}[H]
\centering
\caption{Risultati zero-shot TorchXRayVision.}
\begin{tabular}{llrrr}
\toprule
Dataset & Modello zero-shot & \balacc & \auc & Sens./Spec. \\
\midrule
Kermany & DenseNet-RSNA & 0.668 & 0.825 & 0.944 / 0.393 \\
Kermany & DenseNet-NIH & 0.777 & 0.880 & 0.656 / 0.897 \\
Kermany & DenseNet-All & 0.582 & 0.791 & 0.172 / 0.991 \\
Chittagong & DenseNet-RSNA & 0.813 & 0.905 & 0.805 / 0.821 \\
Chittagong & DenseNet-NIH & 0.635 & 0.705 & 0.531 / 0.739 \\
Chittagong & DenseNet-All & 0.574 & 0.576 & 0.164 / 0.984 \\
\bottomrule
\end{tabular}
\end{table}

Il risultato piu' importante e' che su Chittagong il modello pre-addestrato su
RSNA e' nettamente superiore agli altri zero-shot. Questo supporta la scelta di
un training adult-oriented e pneumonia-specifico. Su Kermany invece il modello
NIH ha la \balacc{} piu' alta: questo conferma che il comportamento dipende dal
dataset e che le conclusioni non possono basarsi su un singolo test.

La stessa conclusione emerge confrontando, sullo stesso dataset Chittagong, le
reti addestrate nella fase iniziale su Kermany con le stesse architetture
addestrate successivamente su RSNA. Il passaggio a RSNA migliora la \balacc{}
media da circa 0.636 a circa 0.801, con un incremento medio di circa 0.165. Con
solo cinque architetture, questo confronto va letto come descrittivo: i delta
vanno da +0.082 a +0.211 e non costituiscono un intervallo inferenziale robusto.

\begin{table}[H]
\centering
\caption{Confronto su Chittagong: training iniziale su Kermany vs training su RSNA.}
\begin{tabular}{lrrr}
\toprule
Modello & Kermany$\rightarrow$Chittagong BA & RSNA$\rightarrow$Chittagong BA & Delta \\
\midrule
PneumoniaNet & 0.634 & 0.830 & +0.196 \\
ResNet18 & 0.608 & 0.819 & +0.211 \\
MobileNetV3 & 0.618 & 0.797 & +0.179 \\
EfficientNet-B0 & 0.636 & 0.795 & +0.159 \\
DenseNet121 & 0.682 & 0.764 & +0.082 \\
\midrule
Media & 0.636 & 0.801 & +0.165 \\
\bottomrule
\end{tabular}
\end{table}

Questo e' il motivo empirico principale per non usare Kermany come dataset di
training finale: a parita' di architettura, il training su RSNA trasferisce
meglio su Chittagong.

\section{Esperimenti post-training su RSNA}

Tutte le reti sono state addestrate sullo stesso subset RSNA bilanciato, con
early stopping e soglia decisionale 0.5. La tabella seguente riporta i risultati
principali con intervalli bootstrap al 95\%.

\small
\begin{longtable}{llrrrlrl}
\caption{Benchmark post-training. BA = balanced accuracy.}\label{tab:benchmark-post-training}\\
\toprule
Dataset & Modello & Sens. & Spec. & BA & BA 95\% CI & \auc & AUC 95\% CI \\
\midrule
\endfirsthead
\toprule
Dataset & Modello & Sens. & Spec. & BA & BA 95\% CI & \auc & AUC 95\% CI \\
\midrule
\endhead
RSNA & PneumoniaNet & 0.918 & 0.874 & 0.896 & [0.877, 0.914] & 0.961 & [0.949, 0.971] \\
RSNA & ResNet18 & 0.956 & 0.926 & 0.941 & [0.926, 0.956] & 0.989 & [0.983, 0.993] \\
RSNA & MobileNetV3 & 0.942 & 0.946 & 0.944 & [0.929, 0.956] & 0.990 & [0.985, 0.993] \\
RSNA & EfficientNet-B0 & 0.944 & 0.956 & \textbf{0.950} & [0.935, 0.962] & 0.989 & [0.983, 0.993] \\
RSNA & DenseNet121 & 0.936 & 0.948 & 0.942 & [0.927, 0.956] & 0.987 & [0.981, 0.992] \\
\midrule
Kermany & PneumoniaNet & 0.944 & 0.432 & 0.688 & [0.654, 0.724] & 0.789 & [0.750, 0.825] \\
Kermany & ResNet18 & 0.956 & 0.487 & 0.722 & [0.685, 0.756] & 0.854 & [0.821, 0.884] \\
Kermany & MobileNetV3 & 0.895 & 0.679 & 0.787 & [0.754, 0.820] & 0.874 & [0.845, 0.901] \\
Kermany & EfficientNet-B0 & 0.867 & 0.769 & 0.818 & [0.786, 0.847] & 0.903 & [0.877, 0.925] \\
Kermany & DenseNet121 & 0.849 & 0.868 & \textbf{0.858} & [0.831, 0.884] & \textbf{0.925} & [0.903, 0.944] \\
\midrule
Chittagong & PneumoniaNet & 0.875 & 0.786 & \textbf{0.830} & [0.795, 0.864] & 0.887 & [0.854, 0.915] \\
Chittagong & ResNet18 & 0.816 & 0.821 & 0.819 & [0.784, 0.852] & 0.893 & [0.861, 0.921] \\
Chittagong & MobileNetV3 & 0.707 & 0.887 & 0.797 & [0.762, 0.832] & \textbf{0.904} & [0.878, 0.930] \\
Chittagong & EfficientNet-B0 & 0.668 & 0.922 & 0.795 & [0.758, 0.828] & 0.894 & [0.863, 0.922] \\
Chittagong & DenseNet121 & 0.602 & 0.926 & 0.764 & [0.729, 0.797] & 0.844 & [0.807, 0.878] \\
\bottomrule
\end{longtable}
\normalsize

\subsection{Controllo della proiezione AP/PA su RSNA}

La proiezione radiografica e' un possibile confondente perche' puo' cambiare la
geometria dell'immagine e, nel subset RSNA, e' associata anche alla classe. Nel
test RSNA sono presenti 479 immagini AP e 521 PA; il gruppo AP contiene 401
positivi e 78 negativi, mentre il gruppo PA contiene 99 positivi e 422 negativi.
La tabella seguente riporta quindi una stratificazione diagnostica su RSNA test,
non una stima causale isolata dell'effetto della proiezione.

\begin{table}[H]
\centering
\caption{Controllo AP/PA su RSNA test. Le metriche sono ricalcolate dagli stessi score per-campione canonici; \(\Delta\)BA = BA(AP) - BA(PA).}
\label{tab:view-position-rsna}
\begin{tabular}{lrrrrr}
\toprule
Modello & AP BA & PA BA & \(\Delta\)BA & AP AUC & PA AUC \\
\midrule
PneumoniaNet & 0.763 & 0.810 & -0.047 & 0.926 & 0.912 \\
ResNet18 & 0.856 & 0.934 & -0.079 & 0.977 & 0.982 \\
MobileNetV3 & 0.883 & 0.925 & -0.043 & 0.977 & 0.984 \\
EfficientNet-B0 & 0.882 & 0.912 & -0.030 & 0.971 & 0.983 \\
DenseNet121 & 0.890 & 0.905 & -0.015 & 0.981 & 0.969 \\
\bottomrule
\end{tabular}
\end{table}

Il controllo mostra che tutti i modelli hanno \balacc{} piu' bassa nel gruppo
AP rispetto al gruppo PA, con differenze tra -0.015 e -0.079. Tuttavia la forte
associazione tra proiezione e classe rende il risultato descrittivo: nel gruppo
AP la specificity e' stimata su soli 78 negativi, mentre nel gruppo PA la
sensitivity e' stimata su soli 99 positivi. Il confondente e' quindi
quantificabile nel dominio RSNA, ma non eliminabile con questo split. Per
Kermany e Chittagong i manifest locali contengono solo percorso e label, senza
metadata DICOM AP/PA; non e' quindi possibile ripetere la stratificazione sui
test esterni senza recuperare o annotare metadata aggiuntivi.

\subsection{Analisi canonica della soglia}

Per evitare una doppia fonte di verita', le metriche di questa analisi sono
ricalcolate dagli score per-campione salvati in
\texttt{outputs/predictions/\{dataset\}/\{model\}.npz}. La soglia di Youden e'
calcolata sullo stesso test set riportato: e' quindi una soglia
\emph{oracle}, ottimistica e non deployabile, usata solo per distinguere il
potere discriminativo dal punto operativo fisso 0.5.

\begin{table}[H]
\centering
\caption{Confronto del ranking a soglia 0.5 e a soglia Youden oracle. Spearman misura la correlazione tra i due ranking entro ciascun dataset.}
\label{tab:threshold-ranking}
\small
\begin{tabular}{lcll}
\toprule
Dataset & Spearman & Ranking a soglia 0.5 & Ranking a Youden \\
\midrule
RSNA & 1.00 & Eff, Mob, Dense, Res, Pneu & Eff, Mob, Dense, Res, Pneu \\
Kermany & 1.00 & Dense, Eff, Mob, Res, Pneu & Dense, Eff, Mob, Res, Pneu \\
Chittagong & 0.10 & Pneu, Res, Mob, Eff, Dense & Eff, Mob, Pneu, Res, Dense \\
\bottomrule
\end{tabular}
\vspace{0.2em}

\footnotesize{Abbreviazioni: Eff = EfficientNet-B0, Mob = MobileNetV3, Dense = DenseNet121, Res = ResNet18, Pneu = PneumoniaNet.}
\end{table}
\normalsize

Il punto operativo 0.5 non riordina i modelli su RSNA e Kermany, ma cambia
fortemente il ranking su Chittagong. In particolare, EfficientNet-B0 passa dal
quarto al primo posto quando si usa la soglia oracle; PneumoniaNet passa dal
primo al terzo. Questo indica che su Chittagong la \balacc{} a soglia fissa
misura anche la posizione operativa scelta, non solo l'architettura.

\begin{table}[H]
\centering
\caption{Gap Kermany--Chittagong a soglia Youden e gap di AUC. Valori positivi indicano prestazioni maggiori su Kermany; valori negativi indicano prestazioni maggiori su Chittagong.}
\label{tab:threshold-gap}
\begin{tabular}{lrrr}
\toprule
Modello & Gap BA@Youden & Gap AUC & Differenza \\
\midrule
PneumoniaNet & -0.110 & -0.098 & -0.012 \\
ResNet18 & -0.058 & -0.039 & -0.019 \\
MobileNetV3 & -0.026 & -0.030 & 0.003 \\
EfficientNet-B0 & -0.012 & 0.009 & -0.021 \\
DenseNet121 & 0.088 & 0.081 & 0.007 \\
\bottomrule
\end{tabular}
\end{table}

Il gap cross-dataset a soglia Youden segue da vicino il gap di AUC:
\(|\Delta(\mathrm{BA}_{Youden})-\Delta(\mathrm{AUC})| \leq 0.021\) per tutte
le cinque architetture. Con \(n=5\) modelli questo va letto come andamento
osservato, non come inferenza statistica generale. Il segno e' pero' informativo:
quattro modelli su cinque hanno \balacc{} oracle piu' alta su Chittagong che su
Kermany, mentre DenseNet121 e' l'unico con gap positivo e quindi con bias di
dominio opposto.

\begin{figure}[H]
\centering
\includegraphics[width=\textwidth]{outputs/threshold_analysis/roc_curves.png}
\caption{Curve ROC sovrapposte per i cinque modelli. I cerchi indicano il punto a soglia 0.5, le croci il punto Youden oracle calcolato sullo stesso test set.}
\label{fig:threshold-roc}
\end{figure}

\subsection{Calibrazione e temperature scaling}

Per verificare se l'instabilita' del ranking su Chittagong dipenda da
miscalibrazione delle probabilita', sono state calcolate metriche di
calibrazione a partire dagli stessi file per-campione usati come sorgente
canonica. La classe positiva e' sempre \pneumonia{}=1. Sono riportati ECE
equal-width, ECE equal-mass e Brier score; l'ECE equal-mass serve soprattutto
nei casi con score molto concentrati vicino a 0 o 1.

\begin{table}[H]
\centering
\caption{Miscalibrazione prima della calibrazione post-hoc. ECE-W = ECE equal-width; ECE-M = ECE equal-mass.}
\label{tab:calibration-ece}
\small
\begin{tabular}{llrrr}
\toprule
Modello & Dataset & ECE-W & ECE-M & Brier \\
\midrule
PneumoniaNet & RSNA & 0.053 & 0.054 & 0.080 \\
PneumoniaNet & Kermany & 0.125 & 0.125 & 0.188 \\
PneumoniaNet & Chittagong & 0.067 & 0.056 & 0.133 \\
ResNet18 & RSNA & 0.019 & 0.019 & 0.042 \\
ResNet18 & Kermany & 0.155 & 0.147 & 0.180 \\
ResNet18 & Chittagong & 0.095 & 0.093 & 0.141 \\
MobileNetV3 & RSNA & 0.026 & 0.023 & 0.042 \\
MobileNetV3 & Kermany & 0.073 & 0.059 & 0.141 \\
MobileNetV3 & Chittagong & 0.130 & 0.121 & 0.154 \\
EfficientNet-B0 & RSNA & 0.014 & 0.009 & 0.037 \\
EfficientNet-B0 & Kermany & 0.049 & 0.058 & 0.119 \\
EfficientNet-B0 & Chittagong & 0.146 & 0.141 & 0.161 \\
DenseNet121 & RSNA & 0.032 & 0.025 & 0.041 \\
DenseNet121 & Kermany & 0.039 & 0.038 & 0.107 \\
DenseNet121 & Chittagong & 0.125 & 0.120 & 0.176 \\
\bottomrule
\end{tabular}
\normalsize
\end{table}

Il confronto e' positivo solo in parte. Chittagong ha ECE equal-mass maggiore di
RSNA per tutti i cinque modelli:
le differenze Chittagong--RSNA sono +0.002 per PneumoniaNet, +0.075 per
ResNet18, +0.098 per MobileNetV3, +0.132 per EfficientNet-B0 e +0.095 per
DenseNet121. Rispetto a Kermany, invece, Chittagong ha ECE maggiore in tre
modelli su cinque: MobileNetV3, EfficientNet-B0 e DenseNet121. Quindi
Chittagong e' sistematicamente peggiore di RSNA in calibrazione, ma non e'
uniformemente peggiore di Kermany.

Sono stati generati anche reliability diagram per ogni modello. Il Brier score
viene usato come controllo binning-free rispetto all'ECE.

E' stato poi applicato temperature scaling binario: per ogni modello si fitta
una temperatura \(T>0\) sui logit della validation RSNA e la stessa \(T\) viene
applicata ai test. Poiche' la trasformazione e' monotona, la ROC-AUC resta
invariata; l'assert nel codice verifica questa proprieta'.

\begin{table}[H]
\centering
\caption{Temperature scaling fittato su RSNA validation e applicato ai test. Le colonne ECE-M e Brier riportano prima$\rightarrow$dopo.}
\label{tab:temperature-scaling}
\small
\begin{tabular}{llrrr}
\toprule
Modello & Dataset & \(T\) & ECE-M & Brier \\
\midrule
PneumoniaNet & RSNA & 0.924 & 0.054$\rightarrow$0.051 & 0.080$\rightarrow$0.080 \\
PneumoniaNet & Kermany & 0.924 & 0.125$\rightarrow$0.132 & 0.188$\rightarrow$0.190 \\
PneumoniaNet & Chittagong & 0.924 & 0.056$\rightarrow$0.056 & 0.133$\rightarrow$0.134 \\
ResNet18 & RSNA & 1.187 & 0.019$\rightarrow$0.021 & 0.042$\rightarrow$0.042 \\
ResNet18 & Kermany & 1.187 & 0.147$\rightarrow$0.135 & 0.180$\rightarrow$0.173 \\
ResNet18 & Chittagong & 1.187 & 0.093$\rightarrow$0.076 & 0.141$\rightarrow$0.137 \\
MobileNetV3 & RSNA & 1.588 & 0.023$\rightarrow$0.022 & 0.042$\rightarrow$0.041 \\
MobileNetV3 & Kermany & 1.588 & 0.059$\rightarrow$0.050 & 0.141$\rightarrow$0.138 \\
MobileNetV3 & Chittagong & 1.588 & 0.121$\rightarrow$0.090 & 0.154$\rightarrow$0.140 \\
EfficientNet-B0 & RSNA & 1.259 & 0.009$\rightarrow$0.010 & 0.037$\rightarrow$0.037 \\
EfficientNet-B0 & Kermany & 1.259 & 0.058$\rightarrow$0.045 & 0.119$\rightarrow$0.118 \\
EfficientNet-B0 & Chittagong & 1.259 & 0.141$\rightarrow$0.135 & 0.161$\rightarrow$0.156 \\
DenseNet121 & RSNA & 0.851 & 0.025$\rightarrow$0.019 & 0.041$\rightarrow$0.040 \\
DenseNet121 & Kermany & 0.851 & 0.038$\rightarrow$0.049 & 0.107$\rightarrow$0.108 \\
DenseNet121 & Chittagong & 0.851 & 0.120$\rightarrow$0.131 & 0.176$\rightarrow$0.181 \\
\bottomrule
\end{tabular}
\normalsize
\end{table}

Su Chittagong, il temperature scaling migliora ECE equal-mass e Brier per
ResNet18, MobileNetV3 ed EfficientNet-B0; e' quasi neutro per PneumoniaNet e
peggiora DenseNet121. Il peggioramento di DenseNet121 non dipende solo dal
binning ECE, perche' anche il Brier passa da 0.1756 a 0.1809. La calibrazione,
pero', non cambia il ranking a soglia 0.5: con \(T>0\) il segno del logit resta
lo stesso e quindi restano identiche anche le predizioni binarie.

Infine, per simulare una scelta di soglia piu' realistica, la soglia Youden e'
stata calcolata su RSNA validation e applicata a Chittagong. Il confronto con
la soglia Youden oracle su Chittagong e' riportato nella tabella seguente.

\begin{table}[H]
\centering
\caption{Soglia Youden fittata su RSNA validation e applicata a Chittagong. L'ultima colonna e' l'upper bound oracle ottenuto scegliendo la soglia su Chittagong test.}
\label{tab:validation-threshold}
\begin{tabular}{lrrr}
\toprule
Modello & Soglia da RSNA val & BA@val-thr & BA@oracle \\
\midrule
PneumoniaNet & 0.605 & 0.813 & 0.832 \\
ResNet18 & 0.719 & 0.809 & 0.825 \\
MobileNetV3 & 0.530 & 0.793 & 0.834 \\
EfficientNet-B0 & 0.324 & 0.809 & 0.836 \\
DenseNet121 & 0.462 & 0.766 & 0.776 \\
\bottomrule
\end{tabular}
\vspace{0.2em}

\footnotesize{ResNet18 ed EfficientNet-B0 appaiono pari a tre decimali in BA@val-thr, ma sui valori pieni non sono identici: 0.8088445 contro 0.8088293. La Spearman e' calcolata sui valori pieni.}
\end{table}

La Spearman tra il ranking ottenuto con soglia da RSNA validation e il ranking
oracle Youden su Chittagong e' 0.20. Quindi una soglia source-domain
deployabile migliora alcuni punti operativi ma non recupera il ranking oracle
del dominio target. La conclusione metodologica e' che la miscalibrazione
contribuisce all'instabilita' su Chittagong, ma una calibrazione globale fittata
solo su RSNA validation non basta a rimuovere lo shift model-specifico del
dominio esterno.

\begin{table}[H]
\centering
\caption{Rapporto tra \balacc{} e costo computazionale, calcolato sui valori completi non arrotondati dei CSV. Valori piu' alti indicano piu' \balacc{} per GMAC, ma non sostituiscono le metriche cliniche principali.}
\begin{tabular}{lrrr}
\toprule
Modello & RSNA BA/GMAC & Kermany BA/GMAC & Chittagong BA/GMAC \\
\midrule
PneumoniaNet & 1.70 & 1.30 & 1.57 \\
ResNet18 & 0.52 & 0.40 & 0.45 \\
MobileNetV3-Large & 4.38 & 3.66 & 3.70 \\
EfficientNet-B0 & 2.47 & 2.13 & 2.07 \\
DenseNet121 & 0.33 & 0.30 & 0.27 \\
\bottomrule
\end{tabular}
\end{table}

Questa lettura costo/prestazione, usata qui come indicatore interno e non come
metrica clinica, mostra che l'aumento di capacita' non produce
un miglioramento proporzionale all'aumento del costo computazionale. Su RSNA,
EfficientNet-B0 ottiene la \balacc{} piu' alta, ma MobileNetV3-Large ha il
miglior rapporto BA/GMAC. Su Kermany DenseNet121 rimane il modello piu'
accurato, ma e' anche il piu' costoso; EfficientNet-B0 e MobileNetV3-Large
offrono un compromesso piu' favorevole tra prestazioni e costo. Su Chittagong
PneumoniaNet e' competitiva in \balacc{}, mentre MobileNetV3-Large resta la
scelta piu' efficiente se si considera il numero di operazioni.

\section{Studio di scaling controllato}

Per isolare l'effetto della capacita' del modello e' stato aggiunto uno studio
di scaling controllato sullo stesso task binario RSNA. L'obiettivo non e'
scegliere un nuovo modello finale, ma misurare quanta capacita' serve per
avvicinarsi alla top-performance. L'intento del protocollo e' mantenere invariati
dataset, seed, augmentation, soglia 0.5, early stopping e loop di training,
facendo variare la complessita' dell'architettura.

Sono state costruite due curve:

\begin{itemize}
  \item \textbf{PneumoniaNet continua}: la rete custom e' stata resa parametrica
  con un moltiplicatore di larghezza \(w\). I canali base
  \(16,32,64,128\) e la testa fully-connected vengono scalati da \(w\).
  Il caso \(w=1.0\) riproduce esattamente la rete originale e mantiene
  300161 parametri.
  \item \textbf{EfficientNet discreta}: sono state addestrate le varianti
  ufficiali ImageNet-pretrained B0, B1, B2 e B3, senza modificare la struttura
  interna delle architetture.
\end{itemize}

La dimensione del training set e' stata resa parametrica nello script
\path{scripts/run_scaling_study.py}; nel risultato riportato qui e' stato usato
il subset RSNA corrente, cioe' 2500 immagini normali e 2500 immagini con
opacita' nel training set. Questa scelta rende il confronto direttamente
comparabile con gli esperimenti principali.

Nei grafici, il ginocchio di saturazione e' stimato come il primo punto in cui
un incremento almeno doppio dei parametri produce meno di +0.01 di \balacc{}
rispetto al miglior punto precedente con al massimo meta' dei parametri.

Va distinta la confrontabilita' di protocollo dalla confrontabilita' di
checkpoint. Lo script di scaling riaddestra i punti della curva e non ricarica
esplicitamente i checkpoint del benchmark principale. Per PneumoniaNet \(w=1.0\)
la rete e' la stessa del benchmark, ma il best checkpoint cambia: epoca 5 nel
benchmark, epoca 16 nello scaling. Questo spiega la differenza tra 0.896 e
0.914 su RSNA e tra 0.830 e 0.756 su Chittagong. I metadata non indicano
differenze di \texttt{epochs}, \texttt{patience}, learning rate o seed; la
spiegazione piu' prudente e' non-determinismo residuo tra esecuzioni separate.
La curva di validation di PneumoniaNet e' quasi piatta (\balacc{} 0.914--0.924),
quindi piccole perturbazioni possono spostare l'argmax tra epoche. Per questo i
valori assoluti di PneumoniaNet vanno riferiti al checkpoint del benchmark; la
curva di scaling serve soprattutto a leggere l'andamento entro-protocollo.

\begin{table}[H]
\centering
\caption{Scaling continuo di PneumoniaNet. BA = balanced accuracy.}
\small
\begin{tabular}{lrrrrr}
\toprule
Width & Parametri & GMAC & RSNA BA & Kermany BA & Chittagong BA \\
\midrule
0.5 & 75457 & 0.134 & 0.922 & 0.716 & 0.815 \\
1.0 & 300161 & 0.527 & 0.914 & 0.686 & 0.756 \\
1.5 & 674113 & 1.181 & 0.922 & 0.722 & 0.780 \\
2.0 & 1197313 & 2.095 & 0.922 & 0.666 & 0.807 \\
3.0 & 2691457 & 4.704 & \textbf{0.924} & 0.677 & \textbf{0.817} \\
4.0 & 4782593 & 8.353 & 0.923 & \textbf{0.724} & 0.795 \\
\bottomrule
\end{tabular}
\normalsize
\end{table}

\begin{table}[H]
\centering
\caption{Scaling discreto EfficientNet. Le varianti sono inizializzate da pesi ImageNet ufficiali. BA = balanced accuracy.}
\small
\begin{tabular}{lrrrrr}
\toprule
Variante & Parametri & GMAC & RSNA BA & Kermany BA & Chittagong BA \\
\midrule
B0 & 4008829 & 0.385 & \textbf{0.950} & \textbf{0.818} & 0.795 \\
B1 & 6514465 & 0.568 & 0.946 & 0.733 & 0.674 \\
B2 & 7702403 & 0.658 & 0.947 & 0.780 & \textbf{0.805} \\
B3 & 10697769 & 0.961 & 0.936 & 0.805 & 0.750 \\
\bottomrule
\end{tabular}
\normalsize
\end{table}

\begin{figure}[H]
\centering
\includegraphics[width=0.92\linewidth]{outputs/scaling_study/balanced_accuracy_vs_params.png}
\caption{Studio di scaling: \balacc{} rispetto al numero di parametri, con asse x logaritmico. La stella indica PneumoniaNet \(w=1.0\); le croci indicano il ginocchio di saturazione stimato.}
\end{figure}

\begin{figure}[H]
\centering
\includegraphics[width=0.92\linewidth]{outputs/scaling_study/balanced_accuracy_vs_gmac.png}
\caption{Studio di scaling: \balacc{} rispetto al costo computazionale in GMAC.}
\end{figure}

Il risultato va letto in modo descrittivo. Su RSNA interno la curva e' quasi
piatta: PneumoniaNet resta tra 0.914 e 0.924, mentre EfficientNet non migliora
passando da B0 a B3. Sui dataset esterni, invece, le curve oscillano senza un
andamento monotono chiaro. Con un solo seed non e' possibile separare l'effetto
della capacita' dal rumore del training; lo studio indica quindi che la
performance interna e' gia' poco sensibile alla dimensione, mentre una
conclusione sullo scaling esterno richiederebbe run multi-seed e, idealmente,
anche valori diversi di \texttt{--train-size}.

\section{Estensione esplorativa multi-classe RSNA}

Come estensione esplorativa, e' stato preparato un task RSNA a tre classi:
\texttt{normal}, \texttt{lung\_opacity} e
\texttt{not\_normal\_no\_lung\_opacity}. Quest'ultima classe corrisponde alla
categoria RSNA esclusa dal task binario: radiografie non normali ma senza
opacita' polmonare annotata. L'obiettivo non e' costruire un secondo benchmark
completo, ma verificare quanto le feature binarie si trasferiscano a una
formulazione leggermente piu' ricca del problema.

Per rendere l'esperimento rapido e collegato al lavoro binario gia' svolto, ogni
modello e' stato inizializzato dal proprio checkpoint RSNA binario. Sono stati
caricati tutti i layer compatibili e sostituita la testa finale con una nuova
uscita a tre logit. Per le architetture pre-addestrate la backbone e' stata
congelata ed e' stata addestrata la testa multi-classe per una epoca; per
PneumoniaNet e' stato usato lo stesso checkpoint binario come inizializzazione.
Questa analisi e' quindi un linear probe sui checkpoint binari, coerente con
l'idea che le feature apprese possano trasferire in modo non uniforme tra task
correlati \cite{yosinski2014transferable}.

\begin{table}[H]
\centering
\caption{Esperimento esplorativo sul task RSNA multi-classe. I modelli sono inizializzati dai rispettivi checkpoint binari RSNA. BA indica macro recall.}
\begin{tabular}{lrrrr}
\toprule
Modello & Accuracy & BA & Macro precision & Macro F1 \\
\midrule
PneumoniaNet & 0.507 & 0.507 & 0.526 & 0.471 \\
ResNet18 & 0.677 & 0.677 & 0.672 & 0.674 \\
MobileNetV3-Large & \textbf{0.692} & \textbf{0.692} & \textbf{0.677} & \textbf{0.677} \\
EfficientNet-B0 & 0.687 & 0.687 & 0.671 & 0.663 \\
DenseNet121 & 0.687 & 0.687 & 0.671 & 0.670 \\
\bottomrule
\end{tabular}
\end{table}

\begin{table}[H]
\centering
\caption{Recall per classe nel task RSNA multi-classe.}
\begin{tabular}{lrrr}
\toprule
Modello & Normal & Lung opacity & Not normal/no opacity \\
\midrule
PneumoniaNet & 0.924 & 0.370 & 0.228 \\
ResNet18 & 0.822 & 0.732 & \textbf{0.478} \\
MobileNetV3-Large & \textbf{0.900} & 0.764 & 0.412 \\
EfficientNet-B0 & 0.894 & \textbf{0.828} & 0.340 \\
DenseNet121 & 0.894 & 0.788 & 0.380 \\
\bottomrule
\end{tabular}
\end{table}

\begin{table}[H]
\centering
\caption{Confusion matrix multi-classe di MobileNetV3-Large su RSNA test, modello migliore per BA macro. Le righe sono le classi vere, le colonne le predizioni.}
\begin{tabular}{lrrr}
\toprule
Classe vera & Pred. normal & Pred. lung opacity & Pred. not normal \\
\midrule
Normal & 450 & 3 & 47 \\
Lung opacity & 17 & 382 & 101 \\
Not normal/no opacity & 121 & 173 & 206 \\
\bottomrule
\end{tabular}
\end{table}

Il risultato principale e' che il trasferimento dal binario al multi-classe e'
limitato. MobileNetV3-Large ottiene la BA macro piu' alta (0.692), ma nessun
modello separa bene la terza classe con una sola testa: il miglior recall su
\texttt{not\_normal\_no\_lung\_opacity} e' 0.478 con ResNet18. La confusion
matrix di MobileNetV3-Large mostra che molti esempi della terza classe vengono
assorbiti da \texttt{normal} o \texttt{lung\_opacity}. Il risultato suggerisce
che un vero benchmark multi-classe richiede fine-tuning end-to-end e analisi
degli errori dedicata.

\section{Test statistici tra modelli}

Per evitare di dichiarare ``migliore'' un modello solo sulla base di piccole
differenze numeriche, e' stato confrontato il modello migliore di ogni dataset
contro gli altri modelli sullo stesso dataset. McNemar valuta differenze nelle
decisioni binarie corretto/errato; DeLong valuta differenze di AUC.

\begin{table}[H]
\centering
\caption{Confronti paired principali. Le p-value riportate sono corrette con Holm-Bonferroni.}
\begin{tabular}{llllrr}
\toprule
Dataset & Riferimento & Confronto & AUC diff. & McNemar p$_H$ & DeLong p$_H$ \\
\midrule
RSNA & EfficientNet & PneumoniaNet & 0.028 & $7.75e{-9}$ & $8.49e{-9}$ \\
RSNA & EfficientNet & ResNet18 & 0.000 & 0.563 & 1.000 \\
RSNA & EfficientNet & MobileNetV3 & -0.001 & 0.563 & 1.000 \\
RSNA & EfficientNet & DenseNet121 & 0.002 & 0.563 & 1.000 \\
\midrule
Kermany & DenseNet121 & PneumoniaNet & 0.136 & $6.03e{-7}$ & $1.20e{-17}$ \\
Kermany & DenseNet121 & ResNet18 & 0.071 & $2.19e{-4}$ & $4.50e{-10}$ \\
Kermany & DenseNet121 & MobileNetV3 & 0.051 & 0.0296 & $2.40e{-5}$ \\
Kermany & DenseNet121 & EfficientNet & 0.022 & 0.0805 & 0.0233 \\
\midrule
Chittagong & PneumoniaNet & ResNet18 & -0.006 & 0.561 & 0.951 \\
Chittagong & PneumoniaNet & MobileNetV3 & -0.017 & 0.247 & 0.397 \\
Chittagong & PneumoniaNet & EfficientNet & -0.007 & 0.247 & 0.951 \\
Chittagong & PneumoniaNet & DenseNet121 & 0.043 & 0.0107 & 0.0092 \\
\bottomrule
\end{tabular}
\end{table}

Questi test sono il risultato piu' importante della fase finale. Su RSNA
EfficientNet-B0 e' superiore a PneumoniaNet, ma non e' statisticamente
distinguibile in modo netto da ResNet18, MobileNetV3 e DenseNet121. Su Kermany
DenseNet121 e' il modello piu' forte, con differenze significative contro quasi
tutti; contro EfficientNet-B0 la differenza resta piccola e va trattata come
borderline, perche' McNemar non raggiunge 0.05 anche dopo correzione. Su
Chittagong PneumoniaNet ha la \balacc{} piu' alta, ma non e' significativamente
superiore a ResNet18, MobileNetV3 o EfficientNet-B0. E' invece migliore di
DenseNet121.

La conclusione e' quindi prudente: \textbf{non esiste una rete migliore in modo
uniforme su tutti i dataset}. Il ranking dipende dal dominio di test e dalla
metrica considerata.

\section{Interpretazione}

I risultati suggeriscono alcuni punti principali.

\paragraph{1. Il training adult-oriented ha senso.}
Rispetto alla fase iniziale, il passaggio a RSNA rende l'esperimento piu'
coerente con l'obiettivo adult-oriented e riduce il rischio di interpretare male
il fallimento su dataset esterni come limite dell'architettura anziche' come
effetto del dominio di training. Il confronto Kermany$\rightarrow$Chittagong vs
RSNA$\rightarrow$Chittagong mostra un incremento medio di circa 0.165 in
\balacc{} a favore del training su RSNA.

\paragraph{2. La generalizzazione cross-dataset resta difficile.}
La letteratura mostra che modelli CXR possono degradare quando cambiano dataset,
label e protocolli di acquisizione \cite{cohen2020limits}. I risultati ottenuti
seguono questa linea: le prestazioni interne RSNA sono alte per tutti i modelli
principali, ma il ranking cambia su Kermany e Chittagong.

La proiezione AP/PA resta un confondente possibile, non una spiegazione causale
dimostrata. Su RSNA test il confondente e' stato quantificato
(Tabella~\ref{tab:view-position-rsna}): AP e PA non hanno la stessa
composizione di classe e tutti i modelli mostrano \balacc{} piu' bassa nel
gruppo AP. I manifest esterni, invece, non contengono metadata AP/PA e quindi
non permettono lo stesso controllo sistematico. L'analisi a soglia Youden mostra
inoltre che la difficolta' relativa dei due dataset esterni non e' uniforme per
tutte le architetture: quattro modelli su cinque hanno \balacc{} oracle piu'
alta su Chittagong che su Kermany, mentre DenseNet121 inverte il segno. La view
position va quindi trattata come confondente parzialmente quantificato su RSNA,
ma non come meccanismo sufficiente a spiegare il trasferimento esterno.

\paragraph{3. La migliore rete dipende dal dataset.}
EfficientNet-B0 e' la scelta piu' forte su RSNA interno, DenseNet121 e' la piu'
forte su Kermany, mentre Chittagong non mostra un vincitore netto tra i primi
modelli. Questo rende la tesi piu' solida: invece di forzare una conclusione
assoluta, si dimostra che la scelta dell'architettura va valutata insieme al
dominio di test.

\paragraph{4. Soglia e separabilita' spiegano fenomeni diversi.}
L'analisi della soglia ridimensiona la spiegazione piu' semplice del caso
DenseNet121. Se il gap Kermany--Chittagong dipendesse soprattutto dal punto
operativo 0.5, dovrebbe ridursi molto usando la soglia Youden. Invece DenseNet121
passa da un gap di \balacc{} 0.094 a un gap oracle 0.088, vicino al gap di AUC
0.081. Per questo modello, quindi, la differenza tra Kermany e Chittagong
riflette soprattutto separabilita' diversa tra i due domini, non solo una soglia
mal posizionata.

La soglia resta pero' importante per il ranking entro Chittagong. A soglia 0.5,
PneumoniaNet e' prima in \balacc{}; a soglia Youden, EfficientNet-B0 diventa
prima e la correlazione di rango tra i due ordinamenti scende a 0.10
(Tabella~\ref{tab:threshold-ranking}). Il risultato corretto e' quindi doppio:
il ranking su Chittagong e' sensibile al punto operativo, ma il gap
cross-dataset di ciascun modello segue da vicino l'AUC
(Tabella~\ref{tab:threshold-gap}). EfficientNet-B0 e' il caso piu' vicino
all'invarianza di dominio esterno, con gap quasi nullo in entrambe le metriche.

\paragraph{5. La calibrazione spiega solo una parte del problema.}
L'analisi ECE/Brier mostra che Chittagong e' peggio calibrato di RSNA per tutte
le architetture, ma non dimostra che la sola calibrazione risolva il ranking.
Temperature scaling, fittato su RSNA validation e applicato ai test, migliora la
calibrazione su Chittagong per ResNet18, MobileNetV3 ed EfficientNet-B0, ma non
per PneumoniaNet e DenseNet121. Questa osservazione riguarda Chittagong: su
RSNA, ad esempio, PneumoniaNet e DenseNet121 migliorano leggermente dopo
temperature scaling. Nel caso DenseNet121 su Chittagong peggiorano sia ECE
equal-mass (0.120--0.131) sia Brier score (0.1756--0.1809), quindi il problema
non e' solo un artefatto del binning ECE.

La calibrazione tramite temperatura migliora le probabilita', ma non cambia le
decisioni a soglia 0.5. Infatti su Chittagong la Spearman tra ranking a soglia
0.5 e ranking Youden oracle resta 0.10 prima e dopo temperature scaling. Anche
la soglia Youden fittata su RSNA validation e applicata a Chittagong resta
lontana dal ranking oracle (\(\rho=0.20\)). La miscalibrazione contribuisce al
problema, ma una calibrazione globale source-domain non rimuove lo shift
model-specifico del dominio Chittagong.

\paragraph{6. PneumoniaNet resta competitiva su Chittagong.}
PneumoniaNet ha molti meno parametri delle architetture pre-addestrate e potrebbe
aver memorizzato meno scorciatoie specifiche del training set RSNA. Questa e'
solo un'ipotesi interpretativa, ma spiega perche' una rete piu' piccola possa
rimanere nel gruppo di testa su Chittagong pur essendo inferiore nel test
interno RSNA. L'analisi GMAC precisa pero' che la compattezza della rete e'
soprattutto parametrica: MobileNetV3-Large ed EfficientNet-B0 sono piu'
efficienti come operazioni per inferenza.

\paragraph{7. Sanity check con TorchXRayVision.}
Il modello TorchXRayVision DenseNet-RSNA ottiene \balacc{} 0.813 su Chittagong,
mentre le reti addestrate nella pipeline locale su RSNA variano tra 0.764 e
0.830. Questa vicinanza conferma che la pipeline locale produce risultati
coerenti con un riferimento esterno gia' pre-addestrato su radiografie toraciche.

\paragraph{8. La terza classe RSNA cambia la natura del problema.}
Nel task multi-classe, \texttt{No Lung Opacity / Not Normal} non e' una classe
clinicamente omogenea come \texttt{Normal} o \texttt{Lung Opacity}. Nel
protocollo usato qui, i modelli non apprendono da zero il task a tre classi:
usano feature gia' addestrate sul binario e una nuova testa multi-classe. Il
basso recall sulla terza categoria indica quindi soprattutto che feature
ottimizzate per il contrasto opacita'/non-opacita' non modellano bene anomalie
non-polmonari eterogenee. MobileNetV3-Large ottiene il miglior equilibrio macro,
mentre ResNet18 e' il modello che recupera piu' casi della terza classe.
L'estensione e' quindi utile come osservazione metodologica sul transfer dal
binario al multi-classe, non come benchmark definitivo del task multi-classe.

\paragraph{9. Lo scaling separa curva interna piatta e rumore esterno.}
Lo studio di scaling controllato mostra che su RSNA PneumoniaNet resta tra 0.914
e 0.924 di \balacc{} pur passando da 75457 a 4782593 parametri, mentre
EfficientNet-B0 e' gia' il punto migliore della propria famiglia. Questa curva
interna e' piatta entro la variabilita' del training, non un plateau monotono
pulito. Sui dataset esterni, inoltre, le curve oscillano senza trend monotono.
Questo rende il singolo seed una limitazione centrale: per la generalizzazione
esterna il rumore della run sembra comparabile all'effetto della capacita'.

\section{Limitazioni}

I risultati vanno letti tenendo presenti alcune limitazioni metodologiche.

\begin{itemize}
  \item \textbf{View position non completamente controllata}: RSNA contiene
  metadata AP/PA e permette una stratificazione sul test interno, dove AP e PA
  risultano associati alla classe e producono prestazioni diverse. I manifest
  esterni, pero', non contengono la proiezione radiografica; la view position
  resta quindi un confondente di acquisizione non controllabile nella
  generalizzazione cross-dataset senza metadata aggiuntivi.
  \item \textbf{Soglia fissa}: le valutazioni principali usano soglia 0.5. Una
  calibrazione per dominio potrebbe cambiare sensitivity e specificity, ma
  introdurrebbe anche un ulteriore livello di adattamento al dataset esterno.
  L'analisi con soglia Youden riportata qui e' volutamente oracle: la soglia e'
  scelta sullo stesso test set, quindi produce un upper bound diagnostico e non
  un protocollo deployabile.
  \item \textbf{Calibrazione source-domain}: il temperature scaling e la soglia
  Youden deployabile sono fittati su RSNA validation, non sui test esterni. Questo
  evita leakage dai test, ma non garantisce calibrazione ottimale su Kermany o
  Chittagong. In particolare, su Chittagong DenseNet121 peggiora sia in ECE sia
  in Brier dopo temperature scaling; il risultato va quindi letto come evidenza
  che una calibrazione globale source-domain puo' non trasferire a tutti i
  modelli e domini.
  \item \textbf{Ranking a soglia 0.5 dopo temperature scaling}: poiche'
  \(T>0\), temperature scaling e' monotono sui logit e non cambia le decisioni a
  soglia 0.5. Puo' migliorare probabilita', ECE e Brier, ma non puo' modificare
  sensitivity, specificity o \balacc{} a soglia 0.5 senza cambiare anche la
  soglia decisionale.
  \item \textbf{Pattern di soglia su cinque architetture}: il parallelismo tra
  gap BA@Youden e gap AUC e' consistente su tutte le cinque architetture
  testate, ma con \(n=5\) va interpretato come osservazione robusta del
  benchmark, non come inferenza statistica generale.
  \item \textbf{Singolo seed sperimentale}: il subset RSNA e il training sono
  stati eseguiti con un seed controllato, ma non sono state ripetute run
  complete multi-seed per stimare la variabilita' del training. Questa
  limitazione e' particolarmente importante nello studio di scaling: sulle
  valutazioni esterne le oscillazioni tra punti della curva sono paragonabili
  all'effetto della capacita', quindi non vanno lette come trend architetturali
  definitivi.
  \item \textbf{Budget effettivo di training}: il benchmark principale seleziona
  PneumoniaNet a epoca 5, mentre la run di scaling \(w=1.0\) arriva al best
  checkpoint a epoca 16 e ottiene una \balacc{} RSNA piu' alta. Il divario RSNA
  tra EfficientNet-B0 e PneumoniaNet puo' quindi riflettere in parte anche la
  variabilita' del training e dell'early stopping, non solo l'architettura.
  \item \textbf{Subset RSNA bilanciato}: la scelta controlla il numero di esempi
  per classe e rende il confronto tra modelli pulito, ma non usa l'intera
  informazione disponibile in RSNA.
  \item \textbf{Scaling controllato su una sola dimensione dati}: lo studio di
  scaling usa il subset RSNA corrente come default. I modelli piu' grandi
  potrebbero essere limitati anche dal numero di immagini disponibili, non solo
  dalla propria capacita'. Per questo lo script espone \texttt{--train-size} e
  \texttt{--seed}, cosi' lo stesso protocollo puo' essere rilanciato con piu'
  esempi e seed diversi.
  \item \textbf{Delta crossover}: il range riportato per il miglioramento da
  training Kermany a training RSNA e' descrittivo e usa solo cinque architetture
  come unita' di confronto. Non e' un intervallo inferenziale sui singoli
  pazienti, perche' per la fase Kermany-trained sono disponibili metriche
  aggregate e non tutte le predizioni campione per campione.
  \item \textbf{Esperimento multi-classe preliminare}: l'estensione a tre classi
  usa checkpoint binari RSNA come inizializzazione e una sola epoca di
  adattamento della testa multi-classe. Non e' quindi un fine-tuning completo
  multi-epoca e non sostituisce un benchmark multi-classe definitivo.
\end{itemize}

\section{Stato dell'arte e posizionamento}

La classificazione automatica di patologie su radiografie toraciche si colloca
in una letteratura ampia, ma eterogenea per popolazioni, label e protocolli di
validazione. Il dataset di Kermany et al. \cite{kermany2018} e' uno dei
riferimenti piu' noti per la classificazione di polmonite pediatrica e dimostra
l'utilita' del transfer learning in imaging medico. In questo lavoro viene pero'
usato come test esterno, non come sorgente di training principale, perche'
l'obiettivo sperimentale e' un'impostazione adult-oriented basata su RSNA
Pneumonia Detection Challenge \cite{rsna_kaggle}. ChestX-ray8/14 di Wang et al.
\cite{wang2017nih} rappresenta invece un riferimento complementare: introduce
un grande dataset CXR multi-label, con label estratte da report radiologici, e
motiva la cautela verso dataset ampi ma debolmente annotati.

Il punto centrale rispetto alla letteratura non e' ottenere un singolo numero
alto su uno split interno, ma valutare la robustezza quando cambia il dominio.
Cohen et al. \cite{cohen2020limits} mostrano che i modelli CXR possono
degradare in modo marcato in trasferimento cross-dataset e che metriche
aggregate apparentemente buone possono nascondere disaccordi tra modelli. La
tesi si posiziona direttamente su questo problema: addestra tutte le
architetture sullo stesso subset RSNA, poi confronta test interno, Kermany e
Chittagong, includendo analisi di soglia, calibrazione e test paired. I modelli
TorchXRayVision \cite{torchxrayvision} sono usati come controllo zero-shot
esterno, per verificare che la pipeline locale produca risultati compatibili con
modelli CXR gia' pre-addestrati su dataset radiografici.

La scelta delle architetture combina una CNN custom compatta con reti standard
di letteratura. ResNet \cite{he2016resnet} e DenseNet \cite{huang2017densenet}
sono riferimenti consolidati per l'uso di connessioni residuali e dense nella
propagazione delle feature; EfficientNet \cite{tan2019efficientnet} introduce
lo scaling composto di profondita', larghezza e risoluzione; MobileNetV3
\cite{howard2019mobilenetv3} rappresenta la famiglia efficiente orientata a
costo computazionale ridotto. Questa selezione consente di separare tre aspetti:
prestazione interna, trasferimento esterno ed efficienza. La lettura congiunta
di accuratezza, parametri e operazioni e' coerente con Canziani et al.
\cite{canziani2016analysis}, che mostrano come il confronto tra reti profonde
debba considerare anche costo computazionale e non solo accuratezza.

Rispetto agli studi che si limitano a confrontare accuracy o AUC, questo lavoro
aggiunge due controlli metodologici. Il primo riguarda il punto operativo:
balanced accuracy, sensitivity e specificity sono calcolate a soglia 0.5, ma
vengono confrontate anche con soglie Youden oracle e con soglie trasferite da
RSNA validation. Il secondo riguarda la calibrazione: seguendo Guo et al.
\cite{guo2017calibration}, il temperature scaling viene fittato su validation
RSNA e applicato ai test senza usare label target-domain. Il risultato mostra
che la calibrazione probabilistica puo' migliorare ECE e Brier, ma non elimina
lo shift model-specifico del ranking su Chittagong.

L'analisi statistica segue infine un'impostazione paired: McNemar
\cite{mcnemar1947} viene usato per confrontare le decisioni binarie sugli
stessi campioni, DeLong \cite{delong1988} per confrontare AUC paired e bootstrap
\cite{efron1979bootstrap} per intervalli descrittivi sulle metriche. La
correzione di Holm \cite{holm1979} limita il rischio di sovrainterpretare
confronti multipli entro lo stesso dataset. L'estensione multi-classe viene
letta alla luce di Yosinski et al. \cite{yosinski2014transferable}: riusare
feature apprese sul task binario e adattare una nuova testa permette un primo
test di trasferibilita', ma non sostituisce un training multi-classe completo.

\section{Conclusione operativa}

Il lavoro sperimentale supporta una conclusione non banale: dopo il training su
RSNA adulto, tutte le architetture principali raggiungono buone prestazioni
interne, ma la generalizzazione esterna non produce un vincitore universale. La
tesi puo' quindi essere formulata come studio comparativo sulla robustezza
cross-dataset, non come semplice ricerca della rete con accuracy piu' alta.
L'analisi computazionale aggiunge un secondo asse di lettura: PneumoniaNet e'
molto compatta in memoria, mentre MobileNetV3-Large ed EfficientNet-B0 sono piu'
efficienti in termini di GMAC e rapporto prestazioni/costo. In particolare,
MobileNetV3-Large ha il miglior rapporto BA/GMAC su tutti e tre i dataset pur
non essendo mai il modello piu' accurato in assoluto. La scelta
dell'architettura dipende quindi non solo dal dominio di test, ma anche dalla
metrica di selezione: accuratezza, robustezza esterna o efficienza
computazionale.

Lo studio di scaling controllato rafforza questa lettura, ma in modo piu'
metodologico che prescrittivo. Su RSNA interno la \balacc{} e' quasi piatta
rispetto alla capacita' per entrambe le famiglie, entro la variabilita' osservata
tra run. Sui dataset esterni, invece, le curve non sono monotone e oscillano in
un intervallo comparabile all'effetto atteso della capacita'. Con un solo seed,
quindi, lo studio non dimostra una legge di scaling esterna: mostra che serve
controllare la variabilita' prima di attribuire i cambiamenti di generalizzazione
alla complessita' del modello.

L'analisi canonica dagli score per-campione chiarisce infine il ruolo della
soglia decisionale. Il ranking a soglia 0.5 e a soglia Youden e' identico su
RSNA e Kermany, ma cambia fortemente su Chittagong. Allo stesso tempo, il gap
Kermany--Chittagong a soglia Youden segue il gap di AUC per tutti i modelli:
la soglia spiega quindi una parte del ranking operativo, ma non elimina le
differenze di separabilita' tra domini. DenseNet121 e' l'outlier direzionale:
e' l'unico modello che resta migliore su Kermany che su Chittagong anche a
soglia oracle.

L'analisi di calibrazione aggiunge un ulteriore vincolo interpretativo. Gli
errori ECE/Brier indicano che Chittagong e' piu' problematico di RSNA dal punto
di vista probabilistico; tuttavia una correzione deployabile fittata su RSNA
validation non basta a rendere stabile il ranking su Chittagong. Il temperature
scaling riduce ECE e Brier per ResNet18, MobileNetV3 ed EfficientNet-B0 su
Chittagong, ma non modifica le decisioni a soglia 0.5 e peggiora DenseNet121.
La soglia Youden scelta su RSNA validation porta la Spearman rispetto al ranking
oracle Chittagong solo a 0.20. Quindi il punto operativo e la calibrazione
contano, ma il domain shift residuo e' model-specifico e non viene rimosso da
una singola calibrazione source-domain.

L'estensione esplorativa multi-classe con la terza classe RSNA esclusa finora e'
stata avviata su tutte le architetture principali. Nel protocollo preliminare,
MobileNetV3-Large ottiene la migliore \balacc{} macro (0.692), seguita da
EfficientNet-B0 e DenseNet121 (0.687) e ResNet18 (0.677). La classe
\texttt{No Lung Opacity / Not Normal} resta poco separata nel linear probe: il
miglior recall osservato e' 0.478 con ResNet18. Questo indica che il passaggio
da binario a multi-classe richiede feature addestrate esplicitamente anche sulla
terza classe, non soltanto una nuova testa sopra rappresentazioni binarie.

\begin{thebibliography}{99}

\bibitem{rsna_kaggle}
RSNA Pneumonia Detection Challenge.
\newblock Kaggle competition page.
\newblock \url{https://www.kaggle.com/c/rsna-pneumonia-detection-challenge}

\bibitem{kermany2018}
D. S. Kermany et al.
\newblock Identifying Medical Diagnoses and Treatable Diseases by Image-Based Deep Learning.
\newblock \textit{Cell}, 172(5):1122--1131.e9, 2018.
\newblock \url{https://pubmed.ncbi.nlm.nih.gov/29474911/}

\bibitem{wang2017nih}
X. Wang, Y. Peng, L. Lu, Z. Lu, M. Bagheri, R. M. Summers.
\newblock ChestX-ray8: Hospital-scale Chest X-ray Database and Benchmarks on Weakly-Supervised Classification and Localization of Common Thorax Diseases.
\newblock CVPR 2017.
\newblock \url{https://arxiv.org/abs/1705.02315}

\bibitem{torchxrayvision}
J. P. Cohen et al.
\newblock TorchXRayVision: A library of chest X-ray datasets and models.
\newblock arXiv:2111.00595, 2021.
\newblock \url{https://arxiv.org/abs/2111.00595}

\bibitem{cohen2020limits}
J. P. Cohen, M. Hashir, R. Brooks, H. Bertrand.
\newblock On the limits of cross-domain generalization in automated X-ray prediction.
\newblock MIDL 2020.
\newblock \url{https://arxiv.org/abs/2002.02497}

\bibitem{he2016resnet}
K. He, X. Zhang, S. Ren, J. Sun.
\newblock Deep Residual Learning for Image Recognition.
\newblock CVPR 2016.
\newblock \url{https://arxiv.org/abs/1512.03385}

\bibitem{huang2017densenet}
G. Huang, Z. Liu, L. van der Maaten, K. Q. Weinberger.
\newblock Densely Connected Convolutional Networks.
\newblock CVPR 2017.
\newblock \url{https://arxiv.org/abs/1608.06993}

\bibitem{tan2019efficientnet}
M. Tan, Q. V. Le.
\newblock EfficientNet: Rethinking Model Scaling for Convolutional Neural Networks.
\newblock ICML 2019.
\newblock \url{https://arxiv.org/abs/1905.11946}

\bibitem{howard2019mobilenetv3}
A. Howard et al.
\newblock Searching for MobileNetV3.
\newblock ICCV 2019.
\newblock \url{https://arxiv.org/abs/1905.02244}

\bibitem{canziani2016analysis}
A. Canziani, A. Paszke, E. Culurciello.
\newblock An Analysis of Deep Neural Network Models for Practical Applications.
\newblock arXiv:1605.07678, 2016.
\newblock \url{https://arxiv.org/abs/1605.07678}

\bibitem{yosinski2014transferable}
J. Yosinski, J. Clune, Y. Bengio, H. Lipson.
\newblock How transferable are features in deep neural networks?
\newblock NeurIPS 2014.
\newblock \url{https://arxiv.org/abs/1411.1792}

\bibitem{mcnemar1947}
Q. McNemar.
\newblock Note on the sampling error of the difference between correlated proportions or percentages.
\newblock \textit{Psychometrika}, 12:153--157, 1947.
\newblock \url{https://doi.org/10.1007/BF02295996}

\bibitem{delong1988}
E. R. DeLong, D. M. DeLong, D. L. Clarke-Pearson.
\newblock Comparing the areas under two or more correlated receiver operating characteristic curves: a nonparametric approach.
\newblock \textit{Biometrics}, 44(3):837--845, 1988.
\newblock \url{https://pubmed.ncbi.nlm.nih.gov/3203132/}

\bibitem{efron1979bootstrap}
B. Efron.
\newblock Bootstrap Methods: Another Look at the Jackknife.
\newblock \textit{The Annals of Statistics}, 7(1):1--26, 1979.
\newblock \url{https://doi.org/10.1214/aos/1176344552}

\bibitem{holm1979}
S. Holm.
\newblock A Simple Sequentially Rejective Multiple Test Procedure.
\newblock \textit{Scandinavian Journal of Statistics}, 6(2):65--70, 1979.
\newblock \url{https://www.jstor.org/stable/4615733}

\bibitem{guo2017calibration}
C. Guo, G. Pleiss, Y. Sun, K. Q. Weinberger.
\newblock On Calibration of Modern Neural Networks.
\newblock ICML 2017.
\newblock \url{https://arxiv.org/abs/1706.04599}

\end{thebibliography}

\end{document}
